% T14 — «ОГЭ-стиль»: чёрно-белый официальный, тонкие линии между задачами, акцент сверху.
% Профиль: математика, ОГЭ задание №1-5 (числовая прямая, дроби). 11 задач, 1 кол.
\documentclass[a4paper,11pt]{article}

\usepackage{../../../../Lessons/_templates/preamble_worksheet}
\usepackage{../_shared/worksheet_check}

% --- Палитра: почти чёрный + один акцентный бордовый штрих сверху ---
\definecolor{wcprimary}{HTML}{1A202C}
\definecolor{wcprimaryLight}{HTML}{718096}
\definecolor{wcprimaryBG}{HTML}{F7FAFC}
\definecolor{wcaccent}{HTML}{9B2C2C}

\geometry{a4paper, top=14mm, bottom=14mm, left=18mm, right=18mm}

% --- Заголовок: акцентная бордовая линия + чёрный текст ---
\renewcommand{\WorksheetTitle}[1]{%
  \par\noindent
  {\color{wcaccent}\rule{\linewidth}{2pt}}\par\vspace{2mm}
  \noindent{\fontsize{16pt}{20pt}\selectfont\bfseries\color{wcprimary}#1}\par
  \vspace{1mm}{\color{wcprimary!60}\rule{\linewidth}{0.4pt}}\par\vspace{4mm}}

% --- TaskBox: «№N» жирно, тонкая линия снизу ---
\renewcommand{\TaskBox}[2]{%
  \par\nopagebreak\noindent
  {\bfseries\color{wcprimary}№#1.}\ #2\par
  \vspace{1mm}{\color{wcprimary!40}\rule{\linewidth}{0.3pt}}\par\vspace{2mm}}

\SetAnswerFieldSize{28mm}{7mm}

\begin{document}

\WorksheetTitle{Тренировочные задания ОГЭ. Числа и вычисления}
\par\noindent{\small\itshape\color{wcprimary!80}Часть 1 \quad·\quad Задания с кратким ответом \quad·\quad Класс 9 \quad·\quad T14 \quad·\quad ID: T14-oge15-2026-05-26}\par\vspace{3mm}
\FirstPageHeader

\TaskBox{1}{Найдите значение выражения $\dfrac{3.7}{0.5 \cdot 1.48}$.\par\vspace{0.5mm}\hfill\textbf{Ответ:}~\AnswerField{1}}

\TaskBox{2}{Какое из чисел отмечено на координатной прямой точкой $A$, если $A$ лежит между $1.2$ и $1.4$? Варианты: $\sqrt{1.8}$, $\sqrt{2.4}$, $\sqrt{1.5}$, $\sqrt{0.9}$.\par\vspace{0.5mm}\hfill\textbf{Ответ:}~\AnswerField{2}}

\TaskBox{3}{Найдите значение выражения $(\sqrt{18} - \sqrt{2})^2$.\par\vspace{0.5mm}\hfill\textbf{Ответ:}~\AnswerField{3}}

\TaskBox{4}{Решите уравнение $x^2 + 5x = 0$. Если корней несколько, в ответе запишите больший.\par\vspace{0.5mm}\hfill\textbf{Ответ:}~\AnswerField{4}}

\TaskBox{5}{Установите соответствие между графиками функций и формулами: $y = -2x + 1$, $y = 2x - 1$, $y = -2x - 1$. В ответе укажите формулу графика, проходящего через точку $(0; 1)$.\par\vspace{0.5mm}\hfill\textbf{Ответ:}~\AnswerField{5}}

\TaskBox{6}{Найдите значение выражения $\dfrac{a^7 \cdot a^{-3}}{a^2}$ при $a = 6$.\par\vspace{0.5mm}\hfill\textbf{Ответ:}~\AnswerField{6}}

\TaskBox{7}{Упростите выражение $(x - 3)^2 - x(x - 6)$.\par\vspace{0.5mm}\hfill\textbf{Ответ:}~\AnswerField{7}}

\TaskBox{8}{Найдите корень уравнения $\dfrac{x}{4} - \dfrac{x - 1}{3} = 2$.\par\vspace{0.5mm}\hfill\textbf{Ответ:}~\AnswerField{8}}

\TaskBox{9}{Прогрессия задана формулой $a_n = 5n - 2$. Найдите $a_8$.\par\vspace{0.5mm}\hfill\textbf{Ответ:}~\AnswerField{9}}

\TaskBox{10}{Решите неравенство $\dfrac{x - 5}{x + 2} \leqslant 0$. Запишите количество целых решений.\par\vspace{0.5mm}\hfill\textbf{Ответ:}~\AnswerField{10}}

\TaskBox{11}{В коробке $12$ красных и $8$ синих карандашей. Какова вероятность вытащить синий? Запишите десятичную дробь.\par\vspace{0.5mm}\hfill\textbf{Ответ:}~\AnswerField{11}}

\WorksheetQR{T14-oge15/L1/v1}

\end{document}
